\section{等式与不等式如何塑造可行域}
\label{sec:09-Convex-Optimization/02-Convex-Sets/02}

考虑一个微型生产计划：两种产品 A 和 B，每件 A 耗 2 单位原料、1 单位工时，利润 3；每件 B 耗 1 单位原料、3 单位工时，利润 4。原料库存 100，总工时上限 90，产量不能为负。写成约束：

\[
2x_A+x_B\le100,\qquad x_A+3x_B\le90,\qquad x_A,x_B\ge0.
\]

把满足这三条不等式的 \((x_A,x_B)\) 画在平面上，你得到的是一个多边形区域——四条直线围出来的交集。目标函数 \(3x_A+4x_B\) 的等值线是一族平行直线。沿利润增长方向平移等值线，最后碰到多边形的位置就是最优解。

你检查几个候选点。原点 \((0,0)\) 利润为零。沿两个坐标轴各走一步：\((0,30)\) 利润 120，\((30,0)\) 利润 90。再试试顶点 \((36,18)\)——代入约束恰好两条紧——利润 \(3\cdot36+4\cdot18=180\)。其他顶点都更低。

为什么最优解总是在顶点？这是偶然吗？答案藏在约束的几何形状里。

\subsection{仿射集：保留整条直线，不只有线段}

在深入多面体之前，先看一种比凸集更强的结构。集合 \(C\) 是\textbf{仿射集}，如果对任意 \(x,y\in C\) 和任意实数 \(\theta\)（不限于 \([0,1]\)），

\[
\theta x+(1-\theta)y\in C.
\]

与凸组合相比，\(\theta\) 不再限制在零到一之间——穿过两点的整条直线都必须留在集合内。仿射集可以写成

\[
C=x_0+V,
\]

其中 \(V\) 是线性子空间。仿射集就是``可能不经过原点的子空间''。

线性方程组的解集

\[
\{x:Ax=b\}
\]

若非空，就是一个仿射集。取一个特解 \(x_0\) 满足 \(Ax_0=b\)，全部解为

\[
x=x_0+z,\qquad z\in\nullsp(A).
\]

这同时告诉你在等式约束下可以向哪些方向移动：只能沿 \(A\) 的零空间方向前进。仿射集一定凸——它甚至比凸集更强——但在整个 \(\R^n\) 中没有内部。例如平面中的一条直线，你无法在它里面放一个二维小球。处理这类集合时应该使用\textbf{相对内部}：相对于它自身张成的仿射空间来看，直线上的开线段完全可以是内部。这一细节在 Slater 条件和锥优化的约束品性中很重要。

\subsection{超平面是一条线性等式；半空间是它的一侧}

当 \(a\ne0\) 时，

\[
H=\{x:a^{\trans}x=b\}
\]

是\textbf{超平面}。向量 \(a\) 与超平面正交——若 \(x,y\in H\)，则 \(a^{\trans}(x-y)=b-b=0\)；\(b/\norm{ a}\) 决定它到原点的有向距离。二维中的超平面是一条直线，三维中是一张平面。

把等式放宽为不等式

\[
a^{\trans}x\le b,
\]

就得到一个\textbf{闭半空间}。任取其中两点 \(x,y\)，

\[
a^{\trans}(\theta x+(1-\theta)y) = \theta a^{\trans}x+(1-\theta)a^{\trans}y
\le b,
\]

所以半空间是凸的。严格不等式给出开半空间，同样是凸的。

多个半空间与超平面的交

\[
P=\{x:Ax\le b,\ Cx=d\}
\]

称为\textbf{多面体}。它可以有界，也可以无限延伸；有界多面体常称多胞体。线性规划的可行域正是多面体——就像开头那个微型生产计划里由四条不等式围出来的区域。

\subsection{为什么线性目标常在顶点处达到最优}

线性函数 \(c^{\trans}x\) 的等值面是一族互相平行的超平面。想象把等值面沿 \(-c\) 方向（即目标下降最快的方向）匀速平移。它最后与多面体接触的位置给出最小值。

接触处通常是一个顶点。原因很直观：多面体的每个面都是平的，等值面也是平的——两张平的东西最后的接触点要么在``角''上（顶点），要么恰好完全贴合在某条边或某个面上。

精确地讲：若线性函数在非空多胞体上有最优解，则至少存在一个极端点是最优的。但并非每个最优解都是顶点。当目标等值面与一条边平行时，整条边上的所有点都最优。

这个区分值得记牢。``线性规划总在顶点有唯一解''是一个流传很广的误解。正确结论只保证至少一个极端点能承载最优值，不保证顶点唯一，也不保证可行域有界。

\subsection{分离超平面：把几何排除变成线性证据}

分离定理是凸分析的枢纽之一。设 \(C\) 是非空闭凸集，点 \(z\) 不属于 \(C\)。你总能放置一个超平面，使 \(C\) 完全在一侧，而 \(z\) 严格位于另一侧：

\[
a^{\trans}x\le b<a^{\trans}z
\qquad\text{对所有 }x\in C.
\]

这给出了一个令人满意的结论：凸集外部的点可以被一条线性不等式\textbf{认证}为不可行。你不需要展示整个集合 \(C\) 来证明 \(z\notin C\)——一条不等式就足够了。

证明的几何直觉很简洁。令 \(x^\star\) 是 \(z\) 在 \(C\) 上的最近点（投影），取

\[
a=z-x^\star.
\]

向量 \(a\) 从 \(C\) 内部指向外部的 \(z\)。如果某个 \(x\in C\) 满足 \(a^{\trans}(x-x^\star)>0\)，那么从 \(x^\star\) 朝 \(x\) 走一小步会离 \(z\) 更近：函数 \(\phi(t)=\norm{ z-x^\star-t(x-x^\star)}_2^2\) 在 \(t=0\) 处的导数为 \(-2a^{\trans}(x-x^\star)<0\)。而凸性保证这些小步仍在 \(C\) 内，这与 \(x^\star\) 是最近点矛盾。因此

\[
a^{\trans}(x-x^\star)\le0
\]

对所有 \(x\in C\) 恒成立。取 \(b=a^{\trans}x^\star\)，则超平面 \(a^{\trans}x=b\) 在 \(x^\star\) 处支撑着 \(C\)，而

\[
a^{\trans}z-b=a^{\trans}(z-x^\star)=\norm{ z-x^\star}_2^2>0,
\]

即 \(z\) 严格位于另一侧。

分离定理不是纯几何的装饰品。线性分类器寻找的正是分离超平面；Lagrange 对偶把``一组约束不可同时满足''转化为线性组合证书；最优性条件也可以理解为：目标的下降方向与可行方向的锥被一个支持超平面隔开。这三种看似不同的工具用的是同一个几何画面。

\subsection{等式约束为何必须仿射}

标准的凸优化问题允许凸不等式约束 \(f_i(x)\le0\)，却只允许仿射等式约束 \(h(x)=0\)。这个限制的原因可以直接从集合层面看透。

若 \(h\) 是非仿射的凸函数，它的次水平集 \(\{x:h(x)\le0\}\) 是凸的，但等值集 \(\{x:h(x)=0\}\) 未必凸。拿最简单的例子

\[
x^2=1
\]

来说，解集是 \(\{-1,1\}\)。取这两个解的中点 \(0\)——\(0^2=0\ne1\)，不在解集里。把一个凸函数强制等于常数，等于在凸集合上切一刀，切出来的往往只有弯曲的边界，线段性质就此丢失。

不等式则不同。\(x^2\le1\) 给出区间 \([-1,1]\)，它是凸的——不等式保留下的是``整个一侧''，凸函数的一侧仍然凸。这就是为什么凸建模里约束总写成一侧的不等式，而等式只留给仿射函数：你想保留的是被平面切出的平直截面，而不是被曲面切出的弯曲线段。
